% !TeX root = ../main.tex
% Chapter 2
\chapter{Theoretical Foundations and Comprehensive Literature Survey}
\label{Chapter2}
\lhead{Chapter 2. \emph{Theoretical Foundations}}

\section{Taxonomy of Deep Generative Speech Architectures}
Neural generative speech synthesis models the probability distribution of acoustic representations conditioned on linguistic text and speaker descriptors. Modern deep generative paradigms partition into four foundational mathematical families:

\subsection{Autoregressive Acoustic Models}
Autoregressive (\acrshort{ar}) architectures formulate speech generation as a sequential next-token prediction task. Representing the acoustic feature sequence as $\mathbf{X} = (\mathbf{x}_1, \mathbf{x}_2, \dots, \mathbf{x}_T)$, the joint probability is factored via the chain rule of probability:
\begin{equation}
    p(\mathbf{X} \mid \mathbf{W}) = \prod_{t=1}^T p(\mathbf{x}_t \mid \mathbf{x}_{<t}, \mathbf{W}; \theta)
    \label{eq:ar_factorization}
\end{equation}
where $\mathbf{x}_{<t} = (\mathbf{x}_1, \dots, \mathbf{x}_{t-1})$ represents the historical acoustic context and $\mathbf{W}$ denotes the conditioning text sequence. Architectures such as Tacotron, AudioLM, and XTTS-v2~\cite{casanova2024xtts} leverage causal multi-head self-attention mechanisms~\cite{vaswani2017attention} to parameterize each conditional step. 

While autoregressive models capture expressive prosody effectively, they present practical computational bottlenecks for edge devices:
\begin{enumerate}[label=(\alph*)]
    \item \textbf{Inference Complexity $O(T)$}: Generating $T$ acoustic frames strictly requires $T$ sequential forward passes through the network. On a mobile CPU, where each forward pass consumes 15--30\,ms, generating a 10-second utterance ($T \approx 937$ frames) demands over 15 to 25 seconds of compute ($\text{RTF} > 2.0$).
    \item \textbf{Exposure Bias and Compounding Errors}: Because the model is trained with teacher-forcing but tested with its own predictions, small deviations at step $t$ compound recursively over time, sometimes resulting in phoneme deletions, unnatural pitch inflections, or repetition loops.
    \item \textbf{KV-Cache Memory Overhead}: Retaining historical key-value tensors across deep transformer stacks inflates RAM usage linearly with sequence length, increasing \gls{lmk} termination risks on 4\,GB Android devices.
\end{enumerate}

\subsection{Variational Autoencoders and Normalizing Flows}
Conditional variational autoencoders (VAEs) and finite normalizing flows (such as VITS~\cite{kim2021vits}) model speech by maximizing the variational evidence lower bound (\acrshort{elbo}). By employing finite normalizing flows parameterized by stacks of invertible affine coupling layers, VITS establishes bijective mappings between an isotropic Gaussian latent prior and the target acoustic distribution. While VITS achieves non-autoregressive parallel synthesis ($O(1)$ forward passes) with fast execution speeds, its mathematical structure relies on Monotonic Alignment Search (MAS) and dynamic programming alignment paths. These aligners degrade in zero-shot cross-lingual transfer, as the internal phonetic duration predictors cannot accommodate unseen multilingual scripts without dedicated acoustic phoneme dictionaries.

\subsection{Score-Based Diffusion Generative Models}
Denoising Diffusion Probabilistic Models (\acrshort{ddpm})~\cite{ho2020ddpm} and continuous score-based generative models~\cite{song2020score} synthesize data by reversing a progressive forward noising process. Formally, a forward \gls{sde} corrupts the data distribution $p_{\text{data}}(\mathbf{x})$ into an isotropic Gaussian distribution $\mathcal{N}(\mathbf{0}, \mathbf{I})$ via an It\^o stochastic differential equation:
\begin{equation}
    d\mathbf{x} = \mathbf{f}(\mathbf{x}, t)dt + g(t)d\mathbf{w}
    \label{eq:forward_sde}
\end{equation}
where $\mathbf{w}$ is standard Brownian motion. Generation is achieved by solving the reverse-time \gls{sde}:
\begin{equation}
    d\mathbf{x} = \left[ \mathbf{f}(\mathbf{x}, t) - g(t)^2 \nabla_{\mathbf{x}} \log p_t(\mathbf{x}) \right] dt + g(t)d\bar{\mathbf{w}}
    \label{eq:reverse_sde}
\end{equation}
where $\nabla_{\mathbf{x}} \log p_t(\mathbf{x})$ is the score function, approximated by a deep neural network $\mathbf{s}_\theta(\mathbf{x}, t)$. 

Although diffusion models generate high-quality audio, their stochastic trajectories are inherently curved. Simulating stochastic differential equations requires numerical solvers across 50 to 1,000 sampling steps. Even accelerated diffusion solvers (such as DPM-Solver) require at least 15 to 30 forward evaluations, incurring unviable latency on mobile ARM processors.

\subsection{Continuous Flow Matching Principles}
To avoid the curved trajectories and stochastic overhead of diffusion SDEs, Lipman et al.~\cite{lipman2022flow} formulated \textit{Continuous Flow Matching} (\acrshort{cfm}). CFM synthesizes data along deterministic probability trajectories governed by ordinary differential equations. By constructing linear conditional displacement interpolants, CFM ensures that individual paths between prior noise and clean speech follow straight lines, providing a sound theoretical framework for low-step numerical solving on resource-constrained edge architectures.

\section{Continuous Normalizing Flows and Optimal Transport Flow Matching}
\label{sec:cfm_theory}
Let $\mathbb{R}^d$ denote the data space representing flattened mel-spectrogram tensors. A \gls{cnf} defines a time-dependent push-forward transformation $\psi_t: \mathbb{R}^d \to \mathbb{R}^d$ for $t \in [0, 1]$, driven by a time-dependent vector field $\mathbf{v}_t: \mathbb{R}^d \to \mathbb{R}^d$ via the initial value problem (\acrshort{ode}):
\begin{equation}
    \frac{d}{dt} \psi_t(\mathbf{x}) = \mathbf{v}_t(\psi_t(\mathbf{x})), \quad \psi_0(\mathbf{x}) = \mathbf{x}
    \label{eq:flow_ivp}
\end{equation}
The vector field induces a continuous probability density path $p_t(\mathbf{x})$ satisfying the continuity equation:
\begin{equation}
    \frac{\partial p_t(\mathbf{x})}{\partial t} + \nabla \cdot (p_t(\mathbf{x}) \mathbf{v}_t(\mathbf{x})) = 0
    \label{eq:continuity_equation}
\end{equation}
where $p_0(\mathbf{x}) = \mathcal{N}(\mathbf{x}; \mathbf{0}, \mathbf{I})$ represents the tractable standard Gaussian prior, and $p_1(\mathbf{x}) = p_{\text{data}}(\mathbf{x})$ denotes the empirical target speech distribution.

\subsection{Conditional Flow Matching (CFM)}
Directly regressing the marginal vector field $\mathbf{v}_t(\mathbf{x})$ is computationally intractable because the marginal density $p_t(\mathbf{x}) = \int p_t(\mathbf{x} \mid \mathbf{x}_1) q(\mathbf{x}_1) d\mathbf{x}_1$ involves an integral over the entire data manifold. Lipman et al.~\cite{lipman2022flow} proved that regressing conditional vector fields $\mathbf{v}_t(\mathbf{x} \mid \mathbf{x}_1)$ conditioned on individual data samples $\mathbf{x}_1$ yields identical parameter gradients:
\begin{theorem}[Equivalence of CFM and Marginal Objectives~\cite{lipman2022flow}]
Let $p_t(\mathbf{x} \mid \mathbf{x}_1)$ be a conditional probability path generated by conditional vector field $\mathbf{v}_t(\mathbf{x} \mid \mathbf{x}_1)$. Define the marginal vector field $\mathbf{v}_t(\mathbf{x}) = \int \mathbf{v}_t(\mathbf{x} \mid \mathbf{x}_1) \frac{p_t(\mathbf{x} \mid \mathbf{x}_1) q(\mathbf{x}_1)}{p_t(\mathbf{x})} d\mathbf{x}_1$. Then the objective:
\begin{equation}
    \mathcal{L}_{\text{CFM}}(\theta) = \mathbb{E}_{t \sim U[0, 1], \mathbf{x}_1 \sim q(\mathbf{x}_1), \mathbf{x} \sim p_t(\mathbf{x} \mid \mathbf{x}_1)} \left\| \mathbf{v}_\theta(\mathbf{x}, t) - \mathbf{v}_t(\mathbf{x} \mid \mathbf{x}_1) \right\|^2
    \label{eq:cfm_objective}
\end{equation}
satisfies $\nabla_\theta \mathcal{L}_{\text{CFM}}(\theta) = \nabla_\theta \mathcal{L}_{\text{Marginal}}(\theta)$.
\end{theorem}

\subsection{Optimal Transport (OT) Displacement Interpolant}
Among the family of conditional probability paths connecting $p_0$ to $p_1$, pairing prior noise samples $\mathbf{x}_0 \sim \mathcal{N}(\mathbf{0}, \mathbf{I})$ with target mel-spectrograms $\mathbf{x}_1 \sim q(\mathbf{x}_1)$ via independent couplings yields the linear displacement interpolant:
\begin{equation}
    \psi_t(\mathbf{x}_0) = (1 - (1 - \sigma_{\text{min}})t) \mathbf{x}_0 + t \mathbf{x}_1
    \label{eq:ot_displacement}
\end{equation}
where $\sigma_{\text{min}} \approx 10^{-5}$ is an infinitesimal stability regularizer. Setting $\sigma_{\text{min}} \to 0$ yields the canonical linear trajectory:
\begin{equation}
    \mathbf{x}_t = (1 - t)\mathbf{x}_0 + t \mathbf{x}_1
    \label{eq:linear_trajectory}
\end{equation}
Differentiating \eref{eq:linear_trajectory} with respect to time $t$ gives the target conditional velocity:
\begin{equation}
    \mathbf{u}_t(\mathbf{x}_t \mid \mathbf{x}_0, \mathbf{x}_1) = \frac{d\mathbf{x}_t}{dt} = \mathbf{x}_1 - \mathbf{x}_0
    \label{eq:ot_velocity}
\end{equation}
The target velocity vector for any individual sample pair is constant across $t$. The neural network $v_\theta(\mathbf{x}_t, t, \mathbf{c})$ is trained to predict the displacement vector $(\mathbf{x}_1 - \mathbf{x}_0)$ given the intermediate state $\mathbf{x}_t$, time step $t$, and conditioning context $\mathbf{c}$:
\begin{equation}
    \mathcal{L}_{\text{OT-CFM}}(\theta) = \mathbb{E}_{t, \mathbf{x}_0, \mathbf{x}_1} \left\| \mathbf{v}_\theta(\mathbf{x}_t, t, \mathbf{c}) - (\mathbf{x}_1 - \mathbf{x}_0) \right\|_2^2
    \label{eq:ot_cfm_loss}
\end{equation}
It is important to note that while conditional trajectories for specific pairs $(\mathbf{x}_0, \mathbf{x}_1)$ are straight lines, the marginal velocity field $\mathbf{v}_t(\mathbf{x})$ learned by the neural network integrates over all possible couplings and can retain residual curvature. Consequently, straight conditional paths do not automatically guarantee that generation will succeed in arbitrarily few integration steps without empirical step tuning or trajectory straightening techniques.

\section{Diffusion Transformers (DiT) and Non-Autoregressive Sequence Modeling}
To parameterize the velocity network $\mathbf{v}_\theta(\mathbf{x}_t, t, \mathbf{c})$, recent architectures have adopted \textit{Diffusion Transformers} (\acrshort{dit})~\cite{peebles2023dit,chen2024f5tts} in place of traditional U-Net convolutional backbones.

\subsection{Backbone Microarchitecture}
In IndicF5, the velocity field estimator is structured as a 22-layer Diffusion Transformer with the following specifications:
\begin{itemize}
    \item \textbf{Transformer Depth}: $L = 22$ identical transformer blocks.
    \item \textbf{Model Hidden Dimension}: $D = 1024$.
    \item \textbf{Attention Multi-Head Configuration}: $H = 16$ attention heads, yielding a per-head dimension of $d_h = \frac{D}{H} = 64$.
    \item \textbf{Feed-Forward Multiplier}: $\text{mult}_{\text{FFN}} = 2$, establishing an intermediate MLP expansion dimension of $D_{\text{FFN}} = 2048$.
    \item \textbf{Total Parameter Footprint}: Approximately 335 million parameters.
\end{itemize}

\subsection{Rotary Position Embeddings (RoPE)}
Unlike sinusoidal positional encodings that inject absolute coordinates into the input layer, IndicF5 employs \textit{Rotary Position Embeddings} (\acrshort{rope}) at every self-attention layer. Given a query vector $\mathbf{q}_m$ and key vector $\mathbf{k}_n$ at sequence positions $m$ and $n$, RoPE applies an orthogonal rotation matrix $\mathbf{R}_{\Theta, m}^d$:
\begin{equation}
    \tilde{\mathbf{q}}_m = \mathbf{R}_{\Theta, m}^d \mathbf{q}_m, \quad \tilde{\mathbf{k}}_n = \mathbf{R}_{\Theta, n}^d \mathbf{k}_n
    \label{eq:rope_rotation}
\end{equation}
such that their inner product encodes relative positional displacement:
\begin{equation}
    \langle \tilde{\mathbf{q}}_m, \tilde{\mathbf{k}}_n \rangle = \mathbf{q}_m^T \mathbf{R}_{\Theta, n-m}^d \mathbf{k}_n = g(\mathbf{q}_m, \mathbf{k}_n, m - n)
    \label{eq:rope_relative}
\end{equation}
This property enables the DiT to generalize to variable-length audio sequences during inference without positional coordinate degradation.

\subsection{ConvNeXt-V2 Text Representation and Fused Input Conditioning}
Text tokens are mapped to a 512-dimensional embedding space. Rather than passing text embeddings to separate cross-attention layers, IndicF5 processes the character token sequence through four ConvNeXt-V2 convolutional blocks incorporating depthwise convolutions, layer normalization, inverted bottlenecks, and Global Response Normalization (GRN). This extracts hierarchical local phonological representations. 

The resulting text representation, the noisy mel tensor $\mathbf{x}_t$, and the reference acoustic prefix $\mathbf{x}_{\text{cond}}$ are concatenated along the channel dimension ($[\mathbf{x}_t \,\|\, \mathbf{x}_{\text{cond}} \,\|\, \mathbf{e}_{\text{text}}]$) and linearly projected into the transformer embedding dimension. The entire composite sequence is then processed through bidirectional self-attention blocks, enabling fused, single-stream representation modeling across linguistic and acoustic modalities.

\section{In-Context Speech Inpainting and Acoustic Timbre Transfer}
Modern zero-shot voice cloning architectures frame speaker identity transfer as an in-context audio inpainting problem~\cite{le2023voicebox,e2tts2024}.

Let $\mathbf{x}_{\text{ref}} \in \mathbb{R}^{C \times T_{\text{ref}}}$ denote the mel-spectrogram extracted from a reference audio clip, and let $\mathbf{w}_{\text{ref}}$ represent its transcript. The model constructs a composite conditioning tensor along the temporal dimension:
\begin{equation}
    \mathbf{x}_{\text{cond}} = \left[ \mathbf{x}_{\text{ref}} \;\Big\|\; \mathbf{0}_{C \times T_{\text{tgt}}} \right] \in \mathbb{R}^{C \times (T_{\text{ref}} + T_{\text{tgt}})}
    \label{eq:conditioning_tensor}
\end{equation}
where $\|$ denotes concatenation along the time axis, and $\mathbf{0}_{C \times T_{\text{tgt}}}$ represents zero-padding across the target generation window.

During training, random spans of audio are masked. The model learns to reconstruct the masked acoustic frames conditioned on the unmasked frames and the complete text sequence $[\mathbf{w}_{\text{ref}} \,\|\, \mathbf{w}_{\text{tgt}}]$.

During inference, the reference audio $\mathbf{x}_{\text{ref}}$ is placed in the unmasked prefix region, while the target region is initialized with Gaussian noise $\mathbf{x}_0 \sim \mathcal{N}(\mathbf{0}, \mathbf{I})$. Throughout the numerical ODE integration, the velocity network $v_\theta$ reads the acoustic context in the prefix frames via bidirectional self-attention. The model extracts formant structure, glottal excitation timbre, speaking rate, and room acoustics from $\mathbf{x}_{\text{ref}}$, propagating these characteristics into the synthesized target frames $\mathbf{x}_{\text{tgt}}$ without requiring explicit speaker classification losses or static d-vector bottlenecks.

\section{Neural Vocoding Paradigms: Transposed Convolutions vs. Fourier Synthesis}
Once the mel-spectrogram tensor $\mathbf{x}_{\text{tgt}} \in \mathbb{R}^{C \times T_{\text{tgt}}}$ is synthesized by the flow-matching backbone, a neural vocoder reconstructs the continuous time-domain speech waveform $y(t) \in \mathbb{R}^{N_{\text{samples}}}$.

\subsection{The HiFi-GAN Paradigm and Edge Bottlenecks}
HiFi-GAN~\cite{kong2020hifigan} has served as a widely adopted neural vocoder. It operates entirely in the time domain, employing a multi-receptive field fusion (MRF) module driven by a cascade of transposed convolutions to upsample 100-channel mel frames (hop size 256) to 24\,kHz raw audio.

While HiFi-GAN produces high-quality audio, its design presents substantial computational overhead on mobile ARM processors:
\begin{enumerate}
    \item \textbf{Memory Bandwidth Overhead}: Upsampling audio by factors of $8 \times 8 \times 2 \times 2 = 256$ requires repeatedly allocating and writing large intermediate feature buffers, exceeding mobile L1/L2 caches.
    \item \textbf{Strided Access Patterns}: Multi-branch dilated convolutions induce strided memory access patterns, hindering continuous SIMD vectorization on ARM \gls{neon} units.
\end{enumerate}

\subsection{Vocos: Fourier Magnitude and Phase Synthesis}
To avoid the physical bottlenecks of transposed convolutions, Siuzdak~\cite{siuzdak2023vocos} developed Vocos. Rather than predicting raw time-domain samples through iterative upsampling, Vocos frames audio generation as complex spectral coefficient prediction.

Because the Short-Time Fourier Transform of a real-valued time-domain signal possesses Hermitian symmetry ($X(-f) = X^*(f)$), the full frequency spectrum is completely determined by $N_{\text{fft}}/2 + 1$ non-redundant frequency bins (513 bins for $N_{\text{fft}} = 1024$). The Vocos backbone consists of ConvNeXt blocks operating directly at the mel frame rate without temporal upsampling, predicting two spectral outputs:
\begin{enumerate}
    \item \textbf{Magnitude Head}: Predicts non-negative magnitude spectra $|\mathbf{X}(f, t)| \in \mathbb{R}^{(N_{\text{fft}}/2 + 1) \times T}$.
    \item \textbf{Phase Spectrogram Head}: Predicts unit complex numbers representing spectral phase angles $\phi(f, t) \in [-\pi, \pi]$ via two output channels $(\cos \phi, \sin \phi)$, normalized such that $\cos^2 \phi + \sin^2 \phi = 1$.
\end{enumerate}

The complex Short-Time Fourier Transform (\acrshort{stft}) representation is assembled via:
\begin{equation}
    \mathbf{X}(f, t) = |\mathbf{X}(f, t)| \odot \left( \cos \phi(f, t) + j \sin \phi(f, t) \right)
    \label{eq:complex_stft}
\end{equation}

The time-domain audio waveform $y(t)$ is reconstructed analytically via the Inverse Short-Time Fourier Transform (\acrshort{istft}):
\begin{equation}
    y(n) = \text{iSTFT}(\mathbf{X}(f, t)) = \frac{\sum_m \mathbf{x}_m(n - mH) w(n - mH)}{\sum_m w^2(n - mH)}
    \label{eq:istft_synthesis}
\end{equation}
where $H = 256$ is the hop size and $w(n)$ is a Hanning synthesis window.

Because the Fourier inversion (\eref{eq:istft_synthesis}) is computed using optimized Fast Fourier Transform (\acrshort{simd}) routines (such as ARM NEON FFT or vDSP), Vocos eliminates transposed convolutions entirely. On mobile ARM CPUs, Vocos runs substantially faster than HiFi-GAN while occupying an INT8 model footprint of $\approx 22$\,MB (13.5M parameters, compressed from 54\,MB in FP32).

\section{Numerical Solvers and Trajectory Dynamics}
During inference, sampling from the velocity field $\mathbf{v}_\theta(\mathbf{x}_t, t, \mathbf{c})$ requires solving the initial value problem:
\begin{equation}
    \frac{d\mathbf{x}_t}{dt} = \mathbf{v}_\theta(\mathbf{x}_t, t, \mathbf{c}), \quad \mathbf{x}_0 \sim \mathcal{N}(\mathbf{0}, \mathbf{I}), \quad t \in [0, 1]
    \label{eq:ode_sampling}
\end{equation}

\subsection{Euler Numerical Integration}
Discretizing the time domain into $N$ intervals, the first-order Euler update step is:
\begin{equation}
    \mathbf{x}_{t + \Delta t} = \mathbf{x}_t + \Delta t \cdot \tilde{\mathbf{v}}_\theta(\mathbf{x}_t, t, \mathbf{c})
    \label{eq:euler_update}
\end{equation}
where $\tilde{\mathbf{v}}_\theta$ incorporates Classifier-Free Guidance (\acrshort{cfg})~\cite{ho2022cfg} with guidance weight $\gamma = 2.0$:
\begin{equation}
    \tilde{\mathbf{v}}_\theta(\mathbf{x}_t, t, \mathbf{c}) = \mathbf{v}_\theta(\mathbf{x}_t, t, \mathbf{c}) + \gamma \left( \mathbf{v}_\theta(\mathbf{x}_t, t, \mathbf{c}) - \mathbf{v}_\theta(\mathbf{x}_t, t, \emptyset) \right) = (1 + \gamma) \mathbf{v}_\theta(\mathbf{x}_t, t, \mathbf{c}) - \gamma \mathbf{v}_\theta(\mathbf{x}_t, t, \emptyset)
    \label{eq:cfg_formula}
\end{equation}
Evaluating \eref{eq:cfg_formula} requires evaluating conditioned and unconditioned inputs (which can be batched together into a single $[2, C, T]$ forward pass).

\subsection{Sway Sampling Parameterization}
While uniform time steps $s_k = \frac{k}{N}$ represent a common discretization, the curvature of the learned vector field is non-uniform across the interval $[0, 1]$. To concentrate sampling steps where vector field trajectories require finer resolution, F5-TTS implements Sway Sampling~\cite{chen2024f5tts} parameterized by coefficient $c \in [-1, 1]$:
\begin{equation}
    t(s) = s + c \cdot \left( \cos\left(\frac{\pi}{2} s\right) - 1 + s \right), \quad s \in [0, 1]
    \label{eq:sway_formula}
\end{equation}
When setting the sway coefficient to $c = -1.0$, this transformation simplifies to:
\begin{equation}
    t(s) = 1 - \cos\left(\frac{\pi}{2} s\right)
    \label{eq:sway_simplified}
\end{equation}
Under this schedule, the derivative $\frac{dt}{ds} = \frac{\pi}{2} \sin\left(\frac{\pi}{2} s\right)$ is smallest near $s = 0$, clustering sampling evaluations more densely in the early generation stages where gross acoustic formant structure is established. In practice, sway sampling aids acoustic stability with fewer solver steps, though rigorous perceptual evaluations remain necessary to establish speech quality across diverse step counts.

\section{Trajectory Straightening and Step Distillation for Edge Efficiency}
To minimize iterative numerical evaluations, prevent battery depletion, and mitigate processor thermal throttling during edge execution on mobile SoCs, two acceleration methodologies have been proposed in the literature:

\subsection{Rectified Flow and Reflow (Liu et al., 2022)}
While optimal transport continuous flow matching produces straight trajectories in expectation, individual sample trajectories can cross in high-dimensional space, forcing the neural network to learn an average velocity that introduces curvature. Liu et al.~\cite{liu2022flowstraight} proved that this curvature can be systematically eliminated through recursive Reflow:

\begin{definition}[The Reflow Operation~\cite{liu2022flowstraight}]
Given a 1-rectified flow model $v_{\theta_1}$ trained on empirical pairs $(\mathbf{x}_0, \mathbf{x}_1)$, generate synthetic pairs $(\mathbf{x}_0, \mathbf{z}_1)$ by numerically integrating $\mathbf{x}_0 \sim \mathcal{N}(\mathbf{0}, \mathbf{I})$ to $t=1$ using $v_{\theta_1}$. The 2-rectified flow model $v_{\theta_2}$ is trained on these synthetic straight-line pairs:
\begin{equation}
    \mathcal{L}_{\text{Reflow}}(\theta_2) = \mathbb{E}_{\mathbf{x}_0, \mathbf{z}_1, t} \left\| \mathbf{v}_{\theta_2}((1-t)\mathbf{x}_0 + t\mathbf{z}_1, t) - (\mathbf{z}_1 - \mathbf{x}_0) \right\|^2
    \label{eq:reflow_loss}
\end{equation}
\end{definition}
By training on coupled pairs generated by the flow itself, trajectory intersections are untangled. The resulting vector field trajectories become straight lines, permitting numerical integration in as few as 4 Euler steps without perceptual degradation.

\subsection{FlashSpeech and Consistency Distillation (Ye et al., 2024)}
FlashSpeech~\cite{flashspeech2024} applies latent consistency models to non-autoregressive speech synthesis. Built upon an adversarial consistency training formulation, FlashSpeech trains a student network to predict the final acoustic representation directly while enforcing consistency between predictions evaluated at adjacent time steps along probability trajectories:
\begin{equation}
    f_\theta(\mathbf{x}_{t_1}, t_1) = f_\theta(\mathbf{x}_{t_2}, t_2) \quad \forall \; t_1, t_2 \in [0, 1]
    \label{eq:consistency_property}
\end{equation}
This self-consistency property collapses the multi-step numerical solving loop into one or two forward evaluations. While FlashSpeech was originally formulated for a distinct latent-consistency speech backbone rather than continuous flow matching, applying consistency distillation principles to IndicF5's DiT architecture constitutes an attractive proposed future direction to enable few-step inference on resource-constrained mobile hardware.

\section{Indic Linguistic Phonetics and Acoustic Corpora}
Indian languages belong predominantly to two distinct linguistic families: the Dravidian family (Tamil, Telugu, Kannada, Malayalam) and the Indo-Aryan family (Hindi, Marathi, Gujarati, Bengali, Punjabi, Odia, Assamese). Unlike English, which employs a Latin alphabet requiring non-deterministic grapheme-to-phoneme conversion (e.g., the pronunciation difference of ``ough'' in \textit{through}, \textit{rough}, and \textit{thought}), Indic scripts are strictly phonetic and syllabic.

\subsection{The Akshara Syllabic Structure}
In Indic orthography, the core phonetic unit is the \textit{akshara}. An akshara is an orthographic syllable structured as:
\begin{equation}
    \text{Akshara} = C_0^k V
    \label{eq:akshara_structure}
\end{equation}
representing zero or more consonants ($C$) followed by an inherent vowel ($V$, typically $/a/$). Consonants are modified by dependent vowel diacritics (\textit{matras}), vowel suppressors (\textit{viramas} or \textit{halant}), and nasalization marks (\textit{anusvaras}). When multiple consonants combine without intervening vowels, they form composite consonant conjuncts (\textit{samyuktaksharas}).

\subsection{Acoustic Duration Properties (IndicVoices)}
The IndicVoices speech corpus~\cite{indicvoices2024} comprises 1,417 hours of verified spontaneous and read speech recorded across 22 Indian languages. Acoustic analysis of this corpus demonstrates that an average Indic akshara possesses an acoustic duration spanning:
\begin{equation}
    \Delta t_{\text{akshara}} \approx 150\,\text{ms} \quad \text{to} \quad 220\,\text{ms}
    \label{eq:akshara_duration}
\end{equation}
In digital mel-spectrogram representations utilizing a 24\,kHz sampling rate and a hop size of 256 samples, each mel frame spans:
\begin{equation}
    \Delta t_{\text{frame}} = \frac{256}{24000} \approx 10.667\,\text{ms}
    \label{eq:frame_duration}
\end{equation}
Because an orthographic akshara typically decomposes into an average of 1.6 to 2.2 Unicode code points (base consonant plus vowel diacritics), dividing the composite syllabic duration (\eref{eq:akshara_duration}) by code point count yields an average temporal density of approximately 65 to 75\,ms per character token. At a frame interval of 10.667\,ms, this corresponds to approximately 6.0 to 7.0 mel frames per Unicode character code point (with $k_{\text{mel}} = 6.5$ frames representing $\approx 69.33$\,ms). This relationship provides the empirical basis for the character-level duration heuristic calibrated in \cref{Chapter3}.

\subsection{IndicF5 and Multilingual Adaptation (AI4Bharat, 2025)}
\label{subsec:indicf5_phir_hera}
IndicF5, formulated by Praveen et~al.~\cite{indicf52025} at AI4Bharat (IIT Madras) in their work \textit{``Phir Hera Fairy: An English Fairytaler is a Strong Faker of Fluent Speech in Low-Resource Indian Languages''}, adapts non-autoregressive continuous flow matching to the multilingual domain of the Indian subcontinent.

The authors evaluated whether an English-pretrained flow-matching checkpoint (F5-TTS~\cite{chen2024f5tts}) could generalize to phonologically distinct Indian languages, testing three optimization strategies:
\begin{enumerate}[leftmargin=*]
    \item \textbf{Training from Scratch}: Initializing the 22-layer DiT backbone with randomized weights on Indian corpora. This yielded slow phonological convergence and unstable mel trajectories due to data scarcity in low-resource languages.
    \item \textbf{Joint Multilingual Pre-training}: Fine-tuning simultaneously on paired English and Indic datasets. While stable, this experienced language competition and reduced prosodic expression in regional dialects.
    \item \textbf{Exclusive Indic Fine-Tuning of the English Prior}: Fine-tuning the pre-trained English F5-TTS checkpoint exclusively on Indic speech data. This strategy achieved the highest naturalness (MOS) and speaker similarity, demonstrating that flow-matching vector fields learned from dense English acoustic priors provide effective inductive biases that transfer across phonetic boundaries without catastrophic forgetting.
\end{enumerate}

Trained on 1,417 hours of multi-speaker speech sourced from IndicVoices-R~\cite{indicvoices2024}, Rasa, IndicTTS, and LIMMITS across eleven official languages (Assamese, Bengali, Gujarati, Hindi, Kannada, Malayalam, Marathi, Odia, Punjabi, Tamil, and Telugu), IndicF5 demonstrates cross-lingual zero-shot cloning. A reference prompt in one language (such as Kannada or Odia) can guide synthesis in another language (such as Hindi or Tamil) while maintaining vocal timbre, pitch trajectory, and cadence. Furthermore, by defining a unified character tokenizer spanning 187 Unicode glyphs across Indic scripts and Latin characters, the model eliminates external G2P rule tables by learning character-to-acoustic alignments end-to-end.

\section{Edge Machine Learning Runtimes: Runtime Comparative Analysis}
Deploying deep flow-matching models on mobile devices requires an execution engine capable of parsing ONNX or FlatBuffer graphs, executing quantized integer operators, and offloading tensor operations to hardware accelerators. Three primary frameworks exist in the modern mobile ML ecosystem:

\subsection{Sherpa-ONNX}
Developed by the Next-gen Kaldi team~\cite{sherpaonnx2023}, \texttt{sherpa-onnx} is an embedded C++ speech recognition and synthesis engine with Flutter/Dart bindings. Sherpa-ONNX provides specialized C++ loaders for Zipformer, VITS, and Matcha-TTS architectures. However, Sherpa-ONNX currently lacks native C++ loaders for Diffusion Transformers (DiT) and F5-TTS, preventing direct execution of IndicF5.

\subsection{ONNX Runtime Mobile (C++ / Dart FFI)}
Microsoft's ONNX Runtime Mobile~\cite{onnxruntime2021} provides a generic C++ execution engine supporting standard ONNX graphs. It includes the XNNPACK execution provider for multi-threaded ARM NEON CPU acceleration and optional NNAPI/QNN delegates. However, executing custom multi-step ODE flow-matching loops requires developing custom C++ wrappers, shared libraries (\texttt{.so}), and manual Dart FFI bindings.

\subsection{Google LiteRT (Formerly TensorFlow Lite)}
Google LiteRT~\cite{google_litert2024} is designed for mobile ML deployments, offering:
\begin{itemize}
    \item Direct PyTorch conversion via \texttt{litert-torch}~\cite{litert_torch2024}, avoiding intermediate ONNX operator translation issues.
    \item A unified \texttt{CompiledModel} API supporting automatic accelerator selection:
    \begin{itemize}
        \item Qualcomm Hexagon HTP via Qualcomm QNN~\cite{qualcomm_qnn2024}.
        \item MediaTek APU via NeuroPilot~\cite{mediatek_neuropilot2024}.
        \item Multi-core ARM NEON fallback via XNNPACK.
    \end{itemize}
    \item Zero-copy hardware buffer management (\texttt{AHardwareBuffer}), minimizing memory transfers.
\end{itemize}
The architectural trade-offs across these runtimes are detailed in \cref{Chapter3} and form the core of our future roadmap in \cref{Chapter5}.